\documentclass[../main.tex]{subfiles}
\begin{document}
\section{Results, toolkit, and the known landscape}
\label{sec:results}

This section collects the classical mechanisms the rest of Part I composes, the geometric
core that explains them, and a map of what is known across the common families. The program
is assembly, not invention: each mechanism is imported from the literature, and the
contribution is to route concrete families to it and read off explicit, dimension-controlled
constants.

\subsection{The geometric core: curvature of \texorpdfstring{$U$}{U}}
\label{subsec:curvature-core}

Write $\Pstar\propto e^{-U}$. If $\Hess U(x)\succeq mI$ for all $x$ with $m>0$, then $U$ is
$m$-strongly convex and all three constants are controlled by $1/m$ (Theorem~\ref{thm:bakry-emery}).
Curvature is a restoring force: large $m$ means a strong pull back to equilibrium, small
constants, fast mixing, strong concentration. The three ways this fails organize the whole of
Part I --- flat directions ($m\to0$) give large or infinite constants, heavy tails forbid
LSI and $T_2$ (Section~\ref{sec:tier4}), and separated wells give exponentially large
constants (Section~\ref{sec:tier5}).

\subsection{The imported toolkit}
\label{subsec:toolkit}

\begin{theorem}[Bakry--\'Emery curvature criterion]
\label{thm:bakry-emery}
If $U$ is $m$-strongly convex, $\Hess U(x)\succeq m\,I_D$ for all $x$ with $m>0$, then
$\Pstar\propto e^{-U}$ satisfies $\CP\le\CLS\le 1/m$.
\end{theorem}
\noindent\cite{BakryEmery1985,BakryGentilLedoux2014}. The workhorse: the curvature floor
$m=\inf_x\lmin(\Hess U(x))$ is exactly what the GLM spine reads off the prior
(Section~\ref{sec:tier3}), and it is exact for the Gaussian (Section~\ref{sec:tier1}).

\begin{theorem}[Brascamp--Lieb inequality]
\label{thm:brascamp-lieb}
If $U$ is strictly convex ($\Hess U(x)\succ0$ for all $x$), then for all smooth $f$,
\[
  \Var_{\Pstar}(f)\ \le\ \E_{\Pstar}\,\inner{(\Hess U)^{-1}\nabla f}{\nabla f}.
\]
\end{theorem}
\noindent\cite{BrascampLieb1976}. Sharper than the uniform floor of
Theorem~\ref{thm:bakry-emery} when the curvature varies in space; it is the natural tool for
data-informed sharpening (Section~\ref{sec:agenda}). For linear $f(\theta)=\inner{v}{\theta}$
it gives $\Var_{\Pstar}(\inner{v}{\theta})\le\E_{\Pstar}\inner{(\Hess U)^{-1}v}{v}$.

\begin{theorem}[Holley--Stroock bounded perturbation]
\label{thm:holley-stroock}
Let $\Pstar=Z^{-1}e^{-W}\nu$ with $\osc(W)=\sup W-\inf W\le A<\infty$. Then
\[
  \CP(\Pstar)\ \le\ e^{A}\,\CP(\nu),
  \qquad
  \CLS(\Pstar)\ \le\ e^{A}\,\CLS(\nu).
\]
\end{theorem}
\noindent\cite{HolleyStroock1987}. The estimate is vacuous when $\osc(W)=\infty$, which is
exactly the separated-mixture regime (Section~\ref{sec:tier5}).

\begin{theorem}[Tensorization]
\label{thm:tensorization}
For a product measure, $\CP\big(\bigotimes_i\mu_i\big)=\max_i\CP(\mu_i)$, and likewise for
$\CLS$. More generally a conditional decomposition
$\pi(\dd x,\dd y)=\pi_X(\dd x)\,\pi(\dd y\mid x)$ reduces the constant to the marginal and the
conditional pieces.
\end{theorem}
\noindent\cite{BakryGentilLedoux2014}. This is what makes the constants \emph{dimension-free}
for product and conditionally-product structure: the constant cares about the worst single
factor, $\max_i$, not the sum $\sum_i$, hence not the dimension.

\begin{theorem}[Otto--Villani]
\label{thm:otto-villani}
A log-Sobolev inequality with constant $\CLS$ implies the quadratic transportation-cost
inequality \eqref{eq:tci} with $\CTCI\le\CLS$.
\end{theorem}
\noindent\cite{OttoVillani2000,Villani2008}. Combined with Theorem~\ref{thm:bakry-emery}
this is how every LSI bound upgrades, for free, to a $W_2$ guarantee.

\begin{theorem}[One-dimensional Hardy/Muckenhoupt criterion]
\label{thm:hardy-1d}
For a measure on $\R$ (and for $1$D reductions: radial, projected, or conditional slices)
with density $\rho=e^{-U}$ and median $m$, the Poincar\'e constant satisfies
$\CP\asymp\max(B_+,B_-)$ up to universal factors, where
$B_+=\sup_{x>m}\Pstar([x,\infty))\int_m^x\rho(t)^{-1}\dd t$ and symmetrically for $B_-$. The
log-Sobolev criterion is the same with an extra logarithmic tail factor.
\end{theorem}
\noindent\cite{Muckenhoupt1972,BobkovGotze1999}. The sharp two-sided tool in one dimension;
it makes $1$D posteriors (Gaussian, Laplace, Gamma, Student, mixtures) an exact laboratory.

\begin{corollary}[Benign mixtures have dimension-free LSI]
\label{cor:ccnw-mixture-lsi}
A mixture $\pi=\int P_\theta\,\mu(\dd\theta)$ whose components satisfy a common LSI and have
uniformly bounded pairwise $\chi^2$-divergence,
$\sup_{\theta,\theta'}\chi^2(P_\theta\|P_{\theta'})<\infty$, satisfies a dimension-free
log-Sobolev inequality.
\end{corollary}
\noindent\cite{ChenChewiNilesWeed2021}. The benign counterpart of the separated-mixture
obstruction (Section~\ref{sec:tier5}): the bounded-overlap regime where mode separation does
\emph{not} blow up the constant.

\subsection{Lower bounds keep the upper bounds honest}
\label{subsec:lower-bounds}

\begin{remark}[The linear / mode-indicator lower bound]
\label{rem:rayleigh-lower-only}
For any test function $f$, the Rayleigh quotient gives a \emph{lower} bound
$\CP\ge\Var_{\Pstar}(f)/\E_{\Pstar}\norm{\nabla f}^2$. Linear tests $f(x)=\inner{v}{x}$ yield
$\CP\ge\lmax(\Cov_{\Pstar})$ for every measure (exact for Gaussians,
Section~\ref{sec:tier1}); mode-indicator tests expose bottlenecks (Section~\ref{sec:tier5}).
A test-function search can therefore \emph{falsify} a claimed upper bound but never
\emph{certify} one --- upper bounds come only from a template above or a validated numeric.
\end{remark}

\subsection{The known landscape}
\label{subsec:landscape}

The qualitative map across common one-dimensional / product families. ``Yes'' means a finite
constant under the stated structure; the obstructions of Sections~\ref{sec:tier4}--\ref{sec:tier5}
explain the failures.

\begin{center}
\begin{tabularx}{\textwidth}{@{}p{0.30\textwidth} c c c Y@{}}
\toprule
Distribution & Poincar\'e & LSI & $T_2$ & Governing reason\\
\midrule
Gaussian $N(0,\Sigma)$              & exact & exact & exact & quadratic potential\\
Strongly log-concave $\Hess U\succeq mI$ & yes & yes & yes & uniform curvature\\
Products of the above               & yes & yes & yes & tensorization\\
Uniform on compact domain           & yes & usually & usually & bounded support\\
Exponential / Laplace               & yes & \textbf{no} & \textbf{no} & exponential tails\\
Gamma                               & usually & \textbf{no} & \textbf{no} & exponential tail\\
Student-$t$                         & \textbf{no} / weighted & \textbf{no} & \textbf{no} & polynomial tail\\
Cauchy                              & \textbf{no} & \textbf{no} & \textbf{no} & too heavy-tailed\\
Separated Gaussian mixture          & finite, huge & finite, huge & finite, huge & bottleneck\\
General log-concave                 & finite & not nec. & not nec. & may lack curvature\\
\bottomrule
\end{tabularx}
\end{center}

For Bayesian targets specifically, the state of knowledge is uneven: Gaussian and
linear-regression posteriors are solved exactly (Section~\ref{sec:tier1}); Gaussian-prior
convex GLMs admit a crude but certified prior-only bound (Section~\ref{sec:tier3}); and sharp
data-dependent constants, heavy-tailed hierarchical posteriors, and multimodal / quotient
constants are largely open (Section~\ref{sec:agenda}).
\end{document}
